\subsection{整环上的模}

命题23. 在整环 A 上的模 E 中，非自由元素的集合 T 是 E 的一个子模。

若 x 与 y 均非自由，则存在 A 中两个非零元素 \(\alpha\) , \(\beta\) ，使得 \(\alpha x = 0\) 和 \(\beta y = 0\) 。则由于 A 是整环， \(\alpha\beta \neq 0\) ，且由于 A 是交换的， \(\alpha\beta(\lambda x + \mu y) = 0\) 对 A 中所有 \(\lambda\) 和 \(\mu\) 成立 ，因此 \(\lambda x + \mu y\) 不是自由的。

注记。设 E 是任意交换环 A 上的模。若 x 是 E 的非自由元素，则子模 Ax 中的每个元素都是非自由的。另一方面，若 A 含有零因子，则 E 中两个非自由元素之和可能是自由的；例如，在 Z/6Z（视为自身上的模）中，3 和 4 不是自由的，但 \(3 + 4 = 1\) 是自由的。

命题23引出以下定义：

定义5. 在整环A上的模E中，E的挠子模是由非自由元素（也称为E的挠元素）组成的E的子模。

当 E 等于其挠子模时（即 E 的每个元素都被 A 的一个元素 \(\neq0\) 所零化时），E 被称为挠模。当 E 的挠子模约化为 0（即 E 的每个非零元素都是自由的）时，E 被称为（按语言习惯的滥用）无挠模。

自由 A-模的每个子模（特别地，每个投射 A-模）都是无挠的。Z-模 Q 是无挠的。

命题 24. 设 A 是一个整环。对于每个 A-模 E，令 T(E) 表示 E 的挠子模。设 \(f: \mathbf{E} \to \mathbf{E}'\) 是一个 A-线性映射，E 和 \(\mathbf{E}'\) 都是 A-模。

\begin{enumerate}
\def\labelenumi{(\roman{enumi})}
\item
  \(f(\mathbf{T}(\mathbf{E})) \subset \mathbf{T}(\mathbf{E}')\) .
\item
  若 \(f\) 是单射，则 \(f(\mathbf{T}(\mathbf{E})) = \mathbf{T}(\mathbf{E}') \cap f(\mathbf{E})\) .
\item
  若 \(f\) 是满射且 \(\operatorname{Ker}(f) \subset \mathrm{T}(\mathbf{E})\) ，则 \(f(\mathrm{T}(\mathbf{E})) = \mathrm{T}(\mathbf{E}')\) .
\end{enumerate}

断言 (i) 和 (ii) 是显然的。另一方面，如果 \(f\) 是满射且 \(x' \in \mathrm{T}(\mathrm{E}')\) ，则 \(x' = f(x)\) ，其中 \(x \in \mathbf{E}\) ，且根据假设存在 \(\alpha \neq 0\) 于 \(\mathbf{A}\) 中，使得 \(f(\alpha x) = \alpha x' = 0\) ；由此 \(\alpha x \in \operatorname{Ker}(f)\) ，且由假设存在 \(\beta \neq 0\) 于 \(\mathbf{A}\) 中，使得 \(\beta(\alpha x) = 0\) ；由于 \(\beta \alpha \neq 0\) ，故 \(x \in \mathrm{T}(\mathrm{E})\) .

推论 1. 对于每个 A-模 E，E/T(E) 是无挠的。

若 \(f: \mathbf{E} \to \mathbf{E}'\) 是一个 A-线性映射，令 \(f_{\mathrm{T}}\) 表示映射 \(\mathrm{T}(\mathbf{E}) \to \mathrm{T}(\mathbf{E}')\) ，它的图与 \(f\) 在 \(\mathrm{T}(\mathbf{E})\) 上的限制的图相同。利用此记号：

推论 2. 对于 A-模与 A-线性映射的每个正合序列

\[
0 \longrightarrow \mathrm{E} ^ {\prime} \stackrel {f} {\longrightarrow} \mathrm{E} \stackrel {g} {\longrightarrow} \mathrm{E} ^ {\prime \prime}
\]

该序列

\[
0 \longrightarrow \mathrm{T} (\mathrm{E} ^ {\prime}) \xrightarrow {f _ {\mathrm{T}}} \mathrm{T} (\mathrm{E}) \xrightarrow {g _ {\mathrm{T}}} \mathrm{T} (\mathrm{E} ^ {\prime \prime})
\]

是正合的。

\[
\operatorname{Ker} \left(g _ {\mathrm{T}}\right) = \operatorname{Ker} (g) \cap \mathrm{T} (\mathrm{E}) = f \left(\mathrm{E} ^ {\prime}\right) \cap \mathrm{T} (\mathrm{E}) = f \left(\mathrm{T} \left(\mathrm{E} ^ {\prime}\right)\right) = \operatorname{Im} \left(f _ {\mathrm{T}}\right).
\]

命题 25. 设 A 是一个整环，且 \((\mathbf{E}_{\iota})\) 是一族 A-模；则

\begin{enumerate}
\def\labelenumi{(\arabic{enumi})}
\setcounter{enumi}{36}
\tightlist
\item
\end{enumerate}

\[
\mathrm{T} \left(\bigoplus_ {\iota} \mathrm{E} _ {\iota}\right) = \bigoplus_ {\iota} \mathrm{T} (\mathrm{E} _ {\iota}).
\]

设 \((x_{\iota})\) 是 \(\bigoplus_{\iota} E_{\iota}\) 的一个元素，使得 \(x_{\iota} \in T(E_{\iota})\) 对所有 \(\iota\) 成立；则每一个 \(x_{\iota}\) 都被 A 的一个元素 \(\alpha_{\iota} \neq 0\) 所零化，并且可以假设 \(\alpha_{\iota} = 1\) （当 \(x_{\iota} = 0\) 时）；由于族 \((x_{\iota})\) 具有有限支撑，A 的元素 \(\alpha = \prod_{\iota} \alpha_{\iota}\) 是有定义的，并且 \(\neq 0\) ；它显然零化 \(\bigoplus_{\iota} x_{\iota}\) ，因此 \(\bigoplus_{\iota} T(E_{\iota}) \subset T\left(\bigoplus_{\iota} E_{\iota}\right)\) ；其逆命题是显然的。

若 E 和 F 是两个 A-模，则显然 \(T(E \otimes_{A} F)\) 包含 \(T(E) \otimes_{A} F\) 和 \(E \otimes_{A} T(F)\) 的典范像；但可以给出无挠 A-模 E, F 的例子，使得 \(T(E \otimes_{A} F) \neq 0\) （习题 31）。

注意，挠模的无限积不一定是挠模；例如，在 Z-模 \(\prod_{n=1}^{\infty}\left(\mathbf{Z}/p^{n}\mathbf{Z}\right)\) （p 为整数， \textgreater1）中，所有坐标均为 1 的元素是自由的。

命题 26. 设 A 是一个整环，K 是其分式域，E 是一个 A-模，且 \(\mathbf{E}_{(\mathbf{K})} = \mathbf{K} \otimes_{\mathbf{A}} \mathbf{E}\) 是通过扩张算子环得到的 K 上的向量空间；设 \(\phi\) 表示典范 A-线性映射 \(x \mapsto 1 \otimes x\) ，它把 E 映入 \(\mathbf{E}_{(\mathbf{K})}\) .

\begin{enumerate}
\def\labelenumi{(\roman{enumi})}
\item
  \(\mathbf{E}_{(\mathbf{K})}\) 的每一个元素都具有形式 \(\lambda^{-1}\phi(x)\) ，其中 \(\lambda \in \mathbf{A}\) , \(\lambda \neq 0\) ， \(x \in \mathbf{E}\) .
\item
  \(\phi\) 的核是 E 的挠子模 T(E)。
\item
  \(\mathrm{E}_{(\mathbb{K})}\) 的每一个元素都具有形式 \(z = \sum_{i=1}^{n} \xi_i \phi(x_i)\) ，其中 \(\xi_i \in \mathbb{K}\) ， \(x_{i}\in \mathbf{E};\) 对所有 \(i\) ，存在 \(\alpha_{i}\in \mathbf{A}\) ，使得 \(\alpha_{i}\neq 0\) 且 \(\alpha_{i}\xi_{i}\in \mathbf{A}\) ；若 \(\alpha = \prod_{i = 1}^{n}\alpha_{i},\) 则 \(\alpha \neq 0\) 且 \(\alpha \xi_{i} = \beta_{i}\in \mathbf{A}\) 对所有 \(i\) 成立，因此，在 \(\mathrm{E}_{(\mathbf{K})}\) 中，
\end{enumerate}

\[
z = \alpha^ {- 1} (\alpha z) = \alpha^ {- 1} \sum_ {i = 1} ^ {n} \beta_ {i} \phi (x _ {i}) = \alpha^ {- 1} \phi \left(\sum_ {i = 1} ^ {n} \beta_ {i} x _ {i}\right)
\]

因为 \(\phi\) 是 A-线性的。

\begin{enumerate}
\def\labelenumi{(\roman{enumi})}
\setcounter{enumi}{1}
\tightlist
\item
  若 \(x \neq 0\) 在 \(\mathbf{E}\) 中不是自由的，则存在 \(\alpha \neq 0\) 于 \(\mathbf{A}\) 中，使得 \(\alpha x = 0\) ，从而 \(\alpha \phi(x) = \phi(\alpha x) = 0\) 于 \(\mathrm{E}_{(\mathbb{K})}\) 中，这蕴含 \(\phi(x) = 0\) 。反之，假设对某个 \(x \in \mathbf{E}\) 有 \(1 \otimes x = 0\) 于 \(\mathrm{E}_{(\mathbb{K})}\) 中；我们证明 \(x\) 是 \(\mathbf{E}\) 的挠元素。考虑集合 \(\mathfrak{M}\) ，它由 \(\mathbf{K}\) 的单生成子 A-子模组成；在包含关系下它是一个右有向集，因为任意两个元素 \(\alpha, \beta\) （属于 \(\mathbf{K}\) ）可以写成 \(\alpha = \zeta^{-1}\xi, \beta = \zeta^{-1}\eta\) ，其中 \(\xi, \eta, \zeta\) 属于 \(\mathbf{A}\) 且 \(\zeta \neq 0\) ，因此 \(\mathbf{A}. \alpha \subset \mathbf{A}. \zeta^{-1}\) 且 \(\mathbf{A}. \beta \subset \mathbf{A}. \zeta^{-1}\) 。此外， \(\mathbf{K}\) 是诸模 \(\mathbf{M} \in \mathfrak{M}\) 的并，因此可以被视为由诸模 \(\mathbf{M} \in \mathfrak{M}\) 和典范嵌入所定义的正向系统的正向极限（§ 6，第 2 号，注记）。因此在典范同构的意义下， \(\mathrm{E}_{(\mathbb{K})} = \varinjlim (\mathbf{M} \otimes_{\mathbf{A}} \mathbf{E})\) （§ 6，第 3 号，命题 7），且关系 \(1 \otimes x = 0\) 于 \(\mathrm{E}_{(\mathbb{K})}\) 中蕴含存在一个
\end{enumerate}

\(M \in \mathfrak{M}\) ，使得 \(1 \in M\) 且 \(1 \otimes x = 0\) 于张量积 \(M \otimes_{A} E\) 中（集合论，III，§ 7，第 5 号，引理 1）。还可进一步假设（如有必要，把 \(M\) 换成 \(M' \supset M\) ，即 \(K\) 的一个单生成子模） \(M = A. \gamma^{-1}\) ，其中 \(\gamma \in A\) 且 \(\gamma \neq 0\) 。现在映射 \(\xi \mapsto \gamma\xi\) 是 \(M\) 到 A-模 \(A\) 上的一个同构；另一方面，典范同构 \(A \otimes_{A} E \to E\) （§ 3，第 4 号，命题 4）把 \(\xi \otimes x\) 映到元素 \(\xi x\) （它属于 \(E\) ）；因此存在一个同构 \(M \otimes_{A} E \to E\) ，它把张量积 \(\xi \otimes x\) 映到元素 \((\gamma\xi)x\) （属于 \(E\) ）。假设 \(1 \otimes x = 0\) 于 \(M \otimes_{A} E\) 中因此蕴含 \(\gamma x = 0\) .

注记。设 \(\alpha^{-1}\phi(x)\) , \(\beta^{-1}\phi(y)\) 是 \(\mathbf{E}_{(\mathbb{K})}\) 的两个元素，其中 \(\alpha\in\mathbf{A}\) , \(\beta\in\mathbf{A}\) , \(x\in\mathbf{E}\) , \(y\in\mathbf{E}\) , \(\alpha\beta\neq0\) 。要使 \(\alpha^{-1}\phi(x)=\beta^{-1}\phi(y)\) ，必要且充分的条件是 \(\beta x-\alpha y\) 是 \(\mathbf{E}\) 的一个挠元素，因为这个关系等价于 \(\beta\phi(x)=\alpha\phi(y)\) ，后者也可以写成 \(\phi(\beta x-\alpha y)=0\) .

推论 1. 若 E 是无挠 A-模，则典范映射 \(\phi: \mathbf{E} \to \mathbf{E}_{(\mathbb{K})}\) 是单射。

回顾（§ 5，第 1 号）：对于 E 到 K 上向量空间 F 的每个 A-线性映射，存在且仅存在一个 K-线性映射 \(\tilde{f}\colon\mathbf{E}_{(\mathbb{K})}\to\mathbf{F}\) ，使得 \(f=\tilde{f}\circ\phi\) ；我们将称 \(\tilde{f}\) 与 \(f\) 相关联.

推论 2. 设 \(f\) 是把 \(\mathbf{E}\) 映入一个向量空间 \(\mathbf{F}\) （ \(\mathbf{K}\) 上的）的 A-线性映射；若 \(\operatorname{Ker}(f) \subset \mathrm{T}(\mathbf{E})\) ，则 K-线性映射 \(\bar{f}\) （与 \(f\) 相关联的）是单射。

我们把 \(\operatorname{Ker}(\bar{f})\) 的一个元素写成 \(\lambda^{-1}\phi(x)\) 的形式，其中 \(\lambda \in \mathbf{A}\) , \(\lambda \neq 0\) , \(x \in \mathbf{E}\) ；关系 \(\bar{f}(\lambda^{-1}\phi(x)) = 0\) 等价于 \(\lambda^{-1}\bar{f}(\phi(x)) = 0\) 于 \(\mathbf{F}\) 中，从而等价于 \(f(x) = \bar{f}(\phi(x)) = 0\) 。根据假设，这意味着 \(x \in \mathrm{T}(\mathbf{E})\) ，因此 \(\phi(x) = 0\) ，这证明了该推论。

推论 3. 设 E 是一个 A-模，g 是 E 到 K 上向量空间 F 的一个 A-线性映射，使得 \(g(\mathbf{E})\) 生成 F 且 \(\operatorname{Ker}(g) \subset \mathrm{T}(\mathbf{E})\) 。则与 g 相关联的 K-线性映射 \(\bar{g}\) 是 \(\mathrm{E}_{(\mathbb{K})}\) 到 F 上的一个同构。

\(\bar{g}\) 由推论 2 可知是单射的，而 \(g(\mathbf{E})\) 生成 F 这一假设蕴含 \(\bar{g}\) 是满射。

对于每个 A-模 E，向量空间 \(E_{(K)}\) 称为与 E 相伴。对 E 的每个子集 S，S 在 K 上的秩（或按语言习惯的滥用，简称为 S 的秩）是 S 的典范像 \(\phi(S)\) 在 \(E_{(K)}\) 中的秩，换言之（第 2 号，定义 1），即 \(E_{(K)}\) 中由 \(\phi(S)\) 生成的向量子空间在 K 上的维数.

当 E 是无挠 A-模时，通常把它与其典范像 \(\phi(\mathbf{E})\) （在 \(\mathrm{E}_{(\mathbb{K})}\) 中）等同。按照这一约定，E 的每个生成系都包含 \(\mathrm{E}_{(\mathbb{K})}\) 的一组基（第 1 号，定理 2）。特别地：

推论 4. 每个有限生成的 A-模都是有限秩的。

注意，这个推论的逆命题不一定成立；例如，Q 是一个秩为 1 的 Z-模，但它在 Z 上不是有限生成的。

回忆（§ 5，第 1 号）：对于每个线性映射 \(f: \mathbf{E} \to \mathbf{E}'\) （其中 \(\mathbf{E}\) 和 \(\mathbf{E}'\) 是 A-模）， \(f_{(\mathbf{K})}\) 表示 K-线性映射 \(1_{\mathbf{K}} \otimes f: \mathbf{E}_{(\mathbf{K})} \to \mathbf{E}_{(\mathbf{K})}'\) .

命题 27. 对于每一个正合序列

\[
\mathrm{E} ^ {\prime} \xrightarrow {f} \mathrm{E} \xrightarrow {g} \mathrm{E} ^ {\prime \prime}
\]

的A-线性映射，相应的K-线性映射序列

\[
\mathbf {E} _ {(\mathbf {K})} ^ {\prime} \xrightarrow {f _ {(\mathbf {K})}} \mathbf {E} _ {(\mathbf {K})} \xrightarrow {g _ {(\mathbf {K})}} \mathbf {E} _ {(\mathbf {K})} ^ {\prime \prime}
\]

是正合的。

假设 \(g_{(\mathbb{K})}(\lambda^{-1} \otimes x) = 0\) ，其中 \(\lambda \in \mathbf{A}\) , \(\lambda \neq 0\) , \(x \in \mathbb{E}\) ；这等价于 \(\lambda^{-1} \otimes g(x) = 0\) 于 \(\mathrm{E}_{(\mathbb{K})}'\) 中，从而也等价于

\[
1 \otimes g (x) = \lambda (\lambda^ {- 1} \otimes g (x)) = 0;
\]

由命题 26，存在 A 中的 \(\alpha \neq 0\) ，使得 \(\alpha g(x) = 0\) 于 \(\mathbf{E}^{\prime \prime}\) 中，或者也 \(g(\alpha x) = 0\) 。根据假设，因此存在 \(x' \in \mathbf{E}'\) ，使得 \(\alpha x = f(x')\) ，从而 \(\lambda^{-1} \otimes x = f_{(\mathbf{K})}(\alpha^{-1}\lambda^{-1} \otimes x')\) ，这证明了该命题。

推论 1. 若 \(\mathbf{E}'\) 是 \(\mathbf{E}\) 的一个子模，则 \(\mathbf{E}_{(\mathbf{K})}'\) 被典范地等同于 \(\mathbf{E}_{(\mathbf{K})}\) 的一个向量子空间，且 \((\mathbf{E} / \mathbf{E}')_{(\mathbf{K})}\) 被等同于 \(\mathbf{E}_{(\mathbf{K})} / \mathbf{E}_{(\mathbf{K})}'\) .

只需将命题 27 应用于正合序列

\[
0 \rightarrow \mathrm{E} ^ {\prime} \rightarrow \mathrm{E} \rightarrow \mathrm{E} / \mathrm{E} ^ {\prime} \rightarrow 0.
\]

推论 2. 对于每一个 A-线性映射 \(f: \mathrm{E} \to \mathrm{F}\) , \(\operatorname{Ker}(f_{(\mathbf{K})}) = (\operatorname{Ker}(f))_{(\mathbf{K})}\) , \(\operatorname{Im}(f_{(\mathbf{K})}) = (\operatorname{Im}(f))_{(\mathbf{K})}\) , \(\operatorname{Coker}(f_{(\mathbf{K})}) = (\operatorname{Coker}(f))_{(\mathbf{K})}\) 在典范同构的意义下。特别地， \(f_{(\mathbf{K})}\) 为单射（相应地，满射，相应地，零映射）的必要且充分条件是 \(\operatorname{Ker}(f) \subset \mathrm{T}(\mathrm{E})\) （相应地， \(\operatorname{Coker}(f)\) 是一个挠模，相应地， \(\operatorname{Im}(f) \subset \mathrm{T}(\mathrm{F})\) ).

这由推论 1 和命题 26 得出。

推论 3. 设 E 是一个 A-模，且 \((x_{\lambda})_{\lambda \in L}\) 是 E 的一个元素族。要使 \((x_{\lambda})\) 成为一个自由族，必要且充分的条件是：在向量 K-空间 \(\mathbf{E}_{(\mathbb{K})}\) 中，族 \((1 \otimes x_{\lambda})\) 是自由的。

族 \((x_{\lambda})\) 定义了一个 A-线性映射 \(f: \mathbf{A}^{(L)} \to \mathbf{E}\) ，使得 \(f(e_{\lambda}) = x_{\lambda}\) 对所有 \(\lambda \in \mathbf{L}((e_{\lambda})\) （ \(\mathbf{A}^{(L)})\) 的典范基）成立 ，而说 \((x_{\lambda})\) 是自由的就意味着 \(f\) 是单射。只需将推论 2 应用于 \(f\) ，并注意到 \(\mathbf{A}^{(L)}\) 是无挠的（命题 25）。

